% Generated draft: requires mathematical and source review.
\documentclass[11pt]{article}
\usepackage[margin=25mm]{geometry}
\usepackage{fontspec}
\setmainfont{Latin Modern Roman}
\usepackage{amsmath,amssymb,mathtools,unicode-math}
\setmathfont{Latin Modern Math}
\usepackage{enumitem,needspace,xurl,hyperref}
\setlength{\parindent}{0pt}
\setlength{\parskip}{5pt}
\setlength{\emergencystretch}{3em}
\newcommand{\N}{\mathbb{N}}
\newcommand{\Z}{\mathbb{Z}}
\newcommand{\Q}{\mathbb{Q}}
\newcommand{\R}{\mathbb{R}}
\newcommand{\C}{\mathbb{C}}
\newcommand{\F}{\mathbb{F}}
\newcommand{\A}{\mathbb{A}}
\newcommand{\PP}{\mathbb{P}}
\newcommand{\E}{\mathbb{E}}
\newcommand{\eps}{\varepsilon}
\newcommand{\dd}{\,\mathrm{d}}
\begin{document}
\section*{UnsolvedMath ID: 20000380}
\textbf{Source problem code:} catalogue entry
\subsection{Definitions and mathematical statement}
Melanie Matchett Wood's question asks for an integral version of the BKLPR heuristics, replacing \(\Z_p\) by \(\Z\), bounding the integral data by height \(B\), calibrating this bound using rank-two elliptic curves, and predicting rank three for comparison with analytic random-matrix theory. The source does not define the integral ensemble, height normalization or rank-two calibration input.

The attached aid uses the Park--Poonen--Voight--Wood box model. Let \(\operatorname{Alt}_n(\Z;B_{\mathrm{mat}})\) consist of alternating matrices \(A\), with \(A^{\mathsf T}=-A\), diagonal zero and upper-triangular entries bounded in absolute value by \(B_{\mathrm{mat}}\). Corank is \(n-\operatorname{rank}_{\Q}A\), with the same parity as \(n\). For fixed \(n\) and \(1\leq r\leq n\) with \(r\equiv n\pmod2\), the aid attributes the counting exponent
\[
\#\{A:\operatorname{corank}A\geq r\}
=B_{\mathrm{mat}}^{n(n-r)/2+o(1)}
\]
to the unconditional PPVW lattice theorem. Dividing by the ambient exponent \(n(n-1)/2\) gives corank-tail exponent \(-n(r-1)/2\). The displayed formula does not apply to the unrestricted corank-zero case, and parity-incompatible thresholds require adjustment.

For a height family of elliptic curves of size \(H^{\delta+o(1)}\), let \(\beta_r\) denote its predicted rank-at-least-\(r\) polynomial counting exponent in the applicable growth regime. The aid conditionally assumes the slow-growing-dimension proxy, parity equidistribution and transfer from matrix corank to Mordell--Weil rank. If
\[
B_{\mathrm{mat}}^n=H^{\alpha+o(1)},
\]
its calibration gives \(\beta_2=\delta-\alpha/2\), hence
\[
\beta_r=(r-1)\beta_2-(r-2)\delta,
\qquad\beta_{r+1}-\beta_r=\beta_2-\delta.
\]
This conditional affine dependence means that a reliable nonaffine sequence cannot be corrected merely by retuning \(B_{\mathrm{mat}}\) and \(n\) within the simple ensemble. A fixed-\(n\) lattice theorem alone does not establish the growing-\(n\) proxy or arithmetic transfer.

For all elliptic curves, the aid reports calibration
\[
B_{\mathrm{mat}}^n=H^{1/12+o(1)},\qquad
\#\{\text{rank at least three}\}=H^{3/4+o(1)}
\]
in its height convention. For quadratic twists counted by twist height \(D\), it reports rank-three prediction \(D^{1/2+o(1)}\). These are heuristic arithmetic predictions, not unconditional counting theorems. The source does not specify exact-rank versus tail counts or cumulative versus height-band conventions, and negative formal exponents cannot literally describe a nonzero cumulative integer count.

The aid says comparison with analytic random-matrix predictions additionally requires BSD and unresolved regulator discretization. Its historical RMT alternatives are therefore not reported as direct contradictions. The actual calibration evidence and full analytic comparison are absent.
\subsection{Short English statement}
Build an integer-matrix version of BKLPR, calibrate it from rank two and predict rank three. The attached aid reports an unconditional fixed-dimension lattice count but only conditional arithmetic exponent transfer, giving the stated all-curve and twist predictions while preserving the BSD/regulator qualifications on random-matrix comparisons.
\subsection{Sources}
\begin{itemize}
\item Local UnsolvedMath registry, permanent ID \(20000380\), exact public code \texttt{catalogue entry}; canonical location \texttt{aim-algebraic-number-theory-notes.json}, notes[31].
\item AIM workshop ``Arithmetic statistics over finite fields and function fields,'' Question 2, attributed to Melanie Matchett Wood; supplied URL \url{https://aimath.org/pastworkshops/arithstatffieldproblems.pdf}. Its verification confirms the extracted ring means \(\Z_p\). No year is supplied; the literature field is empty. The PDF was not consulted here.
\item Attached research attempt 1, ``Integral BKLPR matrices and the rank-2-to-rank-3 exponent transfer,'' citing the Park--Poonen--Voight--Wood lattice theorem and model without full publication citations. BSD and analytic RMT are mentioned without precise source references.
\end{itemize}
\end{document}
